\begin{abstract_in_korean}
 
\begin{center}
{\LARGE 웨어러블 심전도 모니터링을 위한 경량 ResUMamba 기반 온디바이스 실시간 부정맥 분류}
\end{center}
 
\begin{flushright}
\parbox[t]{0.9\textwidth}
{
\begin{flushright}
{\bf \txtAuthor}\\
전산학과\\
연세대학교 일반대학원
\end{flushright}
}
\end{flushright}
\vspace{1em}
 
웨어러블 ECG patch는 일상생활에서 간헐적으로 발생하는 부정맥을 장기간 모니터링할 수 있지만, 장시간 기록을 모두 수동으로 판독하기에는 한계가 있으며, 클라우드 기반 분석은 추가적인 통신 및 데이터 처리 경로를 필요로 한다. 본 논문에서는 환자 조건화 경로를 포함하는 기존 ResUMamba를 스마트폰에서 실행하기 위해 모델을 경량화하고 CPU 중심의 배포 방법을 적용하였다. 새로운 분류 구조를 제안하는 대신, ResU path의 표준 합성곱을 depthwise separable convolution으로 대체하고 decoder의 transposed convolution을 파라미터 없는 선형 보간으로 대체하여 ResDSUMamba를 구성하였다. 또한 Mamba path의 selective scan은 PyTorch의 표준 ATen 연산을 이용해 다시 작성하고, 고정된 입력 길이 $T=501$에 대해 순차적으로 펼친 정적 그래프로 구성하여 배포하였다.
 
ResDSUMamba는 1.3225 M개의 파라미터와 331.38 M MAC을 가지며, 기존 ResUMamba의 2.4883 M 및 721.50 M과 비교하여 각각 46.9\%와 54.1\% 감소하였다. HiCardi 데이터셋의 레코드 단위 분할에서 ResDSUMamba의 Macro-$F_1$과 Micro-$F_1$은 각각 44.15\%와 83.29\%로, ResUMamba의 42.28\%와 79.27\%보다 높았으며, Macro-Se는 78.00\%에서 75.72\%로 낮아졌다. Galaxy S21+(Exynos 2100)의 Cortex-X1 단일 스레드 FP32 조건에서 XNNPACK 위임은 ResUMamba의 평균 모델 실행시간을 1073.57 ms에서 48.25 ms로 줄였고, DSConv를 적용한 ResDSUMamba는 36.70 ms(P95 38.90 ms)였다. 실제 patch 연동 스트리밍에서는 309.21 s 동안 413개의 beat가 대기열 누적 없이 처리되었고, 애플리케이션에서 보고된 평균 심박수는 80.2 bpm이었다. 심박당 처리시간은 평균 59.4 ms, P95 70.3 ms였으며, 추론 스레드의 CPU 사용률은 7.9\%, Peak Native Heap은 99.2 MB로 측정되었다. 실기기 세션에는 기준 라벨이 없어 온디바이스 분류 정확도는 평가하지 않았다.
 
\vspace{\stretch{1}} \noindent
\hrulefill\\
{\bf 핵심되는 말 : }
\parbox[t]{0.8\textwidth}
{심전도(ECG), 부정맥 분류, 상태 공간 모델(SSMs), 온디바이스 추론, 모바일 CPU 최적화, 웨어러블 헬스케어}
 
\end{abstract_in_korean}
\clearpage
